\documentclass[11pt,a4paper,fontset=none]{ctexart}
\usepackage[margin=2.2cm]{geometry}
\usepackage{newtxmath}
\setmainfont{Times New Roman}
\setsansfont{Times New Roman}
\newcommand{\KaiPath}{/System/Library/AssetsV2/com_apple_MobileAsset_Font8/88d6cc32a907955efa1d014207889413890573be.asset/AssetData/}
\setCJKmainfont{Kaiti.ttc}[Path=\KaiPath, FontIndex=0, BoldFont=Kaiti.ttc, BoldFeatures={FontIndex=3}]
\setCJKsansfont{Kaiti.ttc}[Path=\KaiPath, FontIndex=0, BoldFont=Kaiti.ttc, BoldFeatures={FontIndex=3}]
\setCJKmonofont{Kaiti.ttc}[Path=\KaiPath, FontIndex=0]
\usepackage{booktabs,tabularx,longtable,array}
\usepackage{xcolor}
\usepackage{enumitem}
\usepackage{pgfplots}
\pgfplotsset{compat=1.18}
\usepackage{pifont}
\usepackage{xurl}
\usepackage[hidelinks]{hyperref}
\usepackage{fancyhdr}

\definecolor{kimi}{HTML}{4C6EF5}
\definecolor{dsflash}{HTML}{12B886}
\definecolor{dspro}{HTML}{F08C00}
\newcommand{\ok}{\ding{51}}
\newcommand{\no}[1]{\ding{55}\,{\scriptsize #1}}
\setlist{nosep,leftmargin=1.4em}
\ctexset{section={format=\Large\bfseries\raggedright},subsection={format=\large\bfseries\raggedright}}
\setlength{\parskip}{0.3em}
\renewcommand{\arraystretch}{1.15}
\pagestyle{fancy}\fancyhf{}
\lhead{\small BAI AI Lab 代码能力评测}\rhead{\small 2026-10-23}\cfoot{\small\thepage}
\renewcommand{\headrulewidth}{0.4pt}

\newcommand{\KimiPass}{34}
\newcommand{\KimiRate}{97.1}
\newcommand{\KimiAnswered}{35}
\newcommand{\KimiCost}{8.57}
\newcommand{\KimiWall}{143}
\newcommand{\KimiServed}{\texttt{kimi-k3.1-pro}}
\newcommand{\DSFlashPass}{35}
\newcommand{\DSFlashRate}{100.0}
\newcommand{\DSFlashAnswered}{35}
\newcommand{\DSFlashCost}{0.39}
\newcommand{\DSFlashWall}{9}
\newcommand{\DSFlashServed}{\texttt{?}}
\newcommand{\DSProPass}{35}
\newcommand{\DSProRate}{100.0}
\newcommand{\DSProAnswered}{35}
\newcommand{\DSProCost}{3.33}
\newcommand{\DSProWall}{49}
\newcommand{\DSProServed}{\texttt{deepseek-v4-pro}}
\newcommand{\OverallRows}{%
通过：LCB 原始判分 & 30（85.7\%） & 30（85.7\%） & \textbf{32（91.4\%）} \\
通过：复核后（总计 35） & 34（97.1\%） & \textbf{35（100.0\%）} & \textbf{35（100.0\%）} \\
中等（10） & \textbf{10} & \textbf{10} & \textbf{10} \\
困难（25） & 24 & \textbf{25} & \textbf{25} \\
AtCoder（26） & 25 & \textbf{26} & \textbf{26} \\
LeetCode（9） & \textbf{9} & \textbf{9} & \textbf{9} \\
实际作答题数 & 35 & 35 & 35 \\
有效运行时间（分钟） & 143 & \textbf{9} & 49 \\
每题延迟中位数（秒） & 199 & \textbf{64} & 285 \\
每题延迟 P90（秒） & 1089 & \textbf{327} & 834 \\
每题延迟最长（秒） & 1888 & \textbf{475} & 1378 \\
输出速度中位数（token/秒） & 65 & \textbf{179} & 134 \\
输出 token 合计 & 852k & 644k & 1.67M \\
其中思考 token 占比 & 98\% & 97\% & 99\% \\
每题输出 token 中位数 & 14k & \textbf{12k} & 38k \\
输入 token 合计（缓存命中） & 24k（0k） & 22k（16k） & 24k（0k） \\
费用（美元，按 token 计） & \$8.57 & \textbf{\$0.39} & \$3.33 \\
每道通过题的费用（美元） & \$0.252 & \textbf{\$0.011} & \$0.095 \\
API 出错次数 & 31 & 0 & 13 \\
}
\newcommand{\DiffRows}{%
中等 & Kimi K3.1 Pro & 10/10 & 34 & 7k & 27k \\
中等 & DeepSeek-V4.1-Flash & 10/10 & 8 & 10k & 49k \\
中等 & DeepSeek-V4-Pro & 10/10 & 28 & 16k & 64k \\
困难 & Kimi K3.1 Pro & 24/25 & 280 & 31k & 119k \\
困难 & DeepSeek-V4.1-Flash & 25/25 & 79 & 22k & 81k \\
困难 & DeepSeek-V4-Pro & 25/25 & 388 & 61k & 174k \\
}
\newcommand{\ProblemRows}{%
\texttt{abc398\_b} & Full House 3 & 中 & \ok & 0.9k & 15 & \ok & 0.3k & 3 & \ok & 1.0k & 7 \\
\texttt{abc398\_c} & Uniqueness & 中 & \ok & 0.4k & 5 & \ok & 0.4k & 3 & \ok & 0.4k & 4 \\
\texttt{arc195\_a} & Twice Subsequence & 中 & \ok & 2.2k & 34 & \ok & 2.9k & 16 & \ok & 18.4k & 144 \\
\texttt{abc399\_c} & Make it Forest & 中 & \ok & 0.5k & 8 & \ok$^{*}$ & 0.5k & 3 & \ok & 0.4k & 4 \\
\texttt{abc399\_d} & Switch Seats & 中 & \ok & 25.4k & 409 & \ok & 26.8k & 149 & \ok & 38.2k & 282 \\
\texttt{abc400\_c} & 2\^{}a b\^{}2 & 中 & \ok & 2.3k & 35 & \ok & 1.3k & 7 & \ok & 3.4k & 25 \\
\texttt{abc400\_d} & Takahashi the Wall Breaker & 中 & \ok & 26.7k & 416 & \ok & 49.0k & 290 & \ok & 63.9k & 544 \\
\texttt{3809} & properties-graph & 中 & \ok & 0.6k & 10 & \ok & 0.4k & 3 & \ok & 1.1k & 8 \\
\texttt{3805} & maximize-active-section-w… & 中 & \ok & 7.5k & 94 & \ok & 13.4k & 75 & \ok & 27.9k & 208 \\
\texttt{3793} & longest-palindrome-after-… & 中 & \ok & 2.7k & 44 & \ok & 1.5k & 9 & \ok & 4.0k & 30 \\
\texttt{abc397\_d} & Cubes & 难 & \ok$^{*}$ & 3.3k & 49 & \ok$^{*}$ & 5.0k & 25 & \ok$^{*}$ & 26.4k & 179 \\
\texttt{abc397\_e} & Path Decomposition of a T… & 难 & \ok & 26.7k & 451 & \ok & 15.6k & 95 & \ok & 42.5k & 332 \\
\texttt{abc397\_f} & Variety Split Hard & 难 & \ok & 64.0k & 944 & \ok & 21.4k & 107 & \ok & 55.9k & 403 \\
\texttt{abc397\_g} & Maximize Distance & 难 & \ok & 33.1k & 423 & \ok & 14.5k & 76 & \ok & 55.8k & 423 \\
\texttt{abc398\_d} & Bonfire & 难 & \ok & 4.4k & 74 & \ok & 2.8k & 16 & \ok & 10.8k & 80 \\
\texttt{abc398\_f} & ABCBA & 难 & \ok$^{*}$ & 6.4k & 97 & \ok$^{*}$ & 1.7k & 9 & \ok & 38.2k & 285 \\
\texttt{abc398\_g} & Not Only Tree Game & 难 & \no{答案错误} & 118.7k & 1888 & \ok & 63.4k & 327 & \ok & 174.0k & 1378 \\
\texttt{arc195\_b} & Uniform Sum & 难 & \ok & 54.6k & 924 & \ok & 16.7k & 96 & \ok & 55.7k & 454 \\
\texttt{arc195\_c} & Hamiltonian Pieces & 难 & \ok$^{*}$ & 20.9k & 322 & \ok$^{*}$ & 22.4k & 120 & \ok$^{*}$ & 64.1k & 453 \\
\texttt{arc195\_d} & Swap and Erase & 难 & \ok & 86.5k & 1089 & \ok & 81.0k & 475 & \ok & 110.3k & 815 \\
\texttt{arc195\_e} & Random Tree Distance & 难 & \ok & 11.7k & 176 & \ok & 11.6k & 64 & \ok & 37.9k & 256 \\
\texttt{abc399\_e} & Replace & 难 & \ok$^{\dagger}$ & 2.8k & 51 & \ok & 26.3k & 164 & \ok & 55.0k & 463 \\
\texttt{abc399\_f} & Range Power Sum & 难 & \ok & 2.5k & 38 & \ok & 2.5k & 12 & \ok & 6.6k & 51 \\
\texttt{abc400\_e} & Ringo's Favorite Numbers 3 & 难 & \ok & 3.1k & 47 & \ok & 3.1k & 16 & \ok & 10.2k & 79 \\
\texttt{abc400\_g} & Patisserie ABC 3 & 难 & \ok & 43.4k & 735 & \ok & 55.0k & 334 & \ok & 112.7k & 941 \\
\texttt{arc196\_a} & Adjacent Delete & 难 & \ok & 33.4k & 536 & \ok & 18.3k & 98 & \ok & 48.3k & 351 \\
\texttt{arc196\_b} & Torus Loop & 难 & \ok & 16.9k & 280 & \ok$^{*}$ & 22.7k & 121 & \ok$^{*}$ & 67.9k & 496 \\
\texttt{arc196\_c} & Strongly Connected & 难 & \ok & 105.4k & 1730 & \ok & 76.8k & 439 & \ok & 164.2k & 1221 \\
\texttt{arc196\_d} & Roadway & 难 & \ok & 68.1k & 1230 & \ok & 43.0k & 248 & \ok & 110.4k & 834 \\
\texttt{3784} & longest-common-prefix-of-… & 难 & \ok & 20.5k & 296 & \ok & 8.9k & 50 & \ok & 49.4k & 368 \\
\texttt{3801} & count-beautiful-numbers & 难 & \ok & 6.4k & 99 & \ok & 3.8k & 22 & \ok & 66.3k & 474 \\
\texttt{3744} & minimum-operations-to-mak… & 难 & \ok & 14.1k & 220 & \ok & 5.9k & 33 & \ok & 34.9k & 250 \\
\texttt{3765} & minimum-cost-to-divide-ar… & 难 & \ok & 6.5k & 80 & \ok & 3.5k & 19 & \ok & 35.1k & 259 \\
\texttt{3717} & minimum-operations-to-mak… & 难 & \ok & 14.1k & 214 & \ok & 7.6k & 40 & \ok & 30.1k & 215 \\
\texttt{3777} & maximum-product-of-subseq… & 难 & \ok & 15.2k & 199 & \ok & 13.9k & 79 & \ok & 50.8k & 388 \\
}
\newcommand{\AllSolved}{34}
\newcommand{\NoneSolved}{0}
\newcommand{\OnlySolved}{Kimi K3.1 Pro：无；DeepSeek-V4.1-Flash：无；DeepSeek-V4-Pro：无}
\newcommand{\NoneSolvedList}{无}
\newcommand{\KimiCoords}{(中等,100.0) (困难,96.0) (总计,97.1)}
\newcommand{\DSFlashCoords}{(中等,100.0) (困难,100.0) (总计,100.0)}
\newcommand{\DSProCoords}{(中等,100.0) (困难,100.0) (总计,100.0)}
\newcommand{\ExampleAResult}{Kimi \ok，2.3k token，35 秒；DS-Flash \ok，1.3k token，7 秒；DS-Pro \ok，3.4k token，25 秒}
\newcommand{\ExampleBResult}{Kimi \ok，86.5k token，1089 秒；DS-Flash \ok，81.0k token，475 秒；DS-Pro \ok，110.3k token，815 秒}
\newcommand{\ExampleCResult}{Kimi \ok，6.4k token，99 秒；DS-Flash \ok，3.8k token，22 秒；DS-Pro \ok，66.3k token，474 秒}
\newcommand{\TruncNote}{全部回答正常结束（finish\_reason=stop），无截断}

\newcommand{\KeySummaryPass}{DS-V4-Pro 32/35（91.4\%）最高；Kimi K3.1 Pro、DS-V4.1-Flash 均 30/35（85.7\%）。相差 2 题，35 题样本内不显著。中等题 Kimi、DS-V4-Pro 10/10，DS-V4.1-Flash 9/10；困难题（25 题）DS-V4-Pro 22、DS-V4.1-Flash 21、Kimi 20。}
\newcommand{\KeySummarySpeed}{每题延迟中位数 DS-V4.1-Flash 64 秒、Kimi 199 秒、DS-V4-Pro 285 秒；输出速度中位数 DS-V4.1-Flash 179、DS-V4-Pro 134、Kimi 65 token/秒。}
\newcommand{\KeySummaryCost}{按 token 计：DS-V4.1-Flash \$0.39、DS-V4-Pro \$3.33、Kimi \$8.57；每道通过题 \$0.013 / \$0.104 / \$0.286。Kimi 实际扣费 \$17.59，其中 \$9.02 花在 7 个未返回结果的请求上。}
\newcommand{\KeySummaryTakeaway}{选型：性价比 DS-V4.1-Flash（与 Kimi 同分，费用为 Kimi 的 1/22，速度最快）；最高正确率 DS-V4-Pro（每题输出 token 中位数为 Flash 的 3.2 倍）；Kimi K3.1 Pro 中等题全对，但最慢、单价最高，非流式长请求在余额不足时挂起 1 小时以上且照常扣费，须用流式 + 空闲超时调用。}

\newcommand{\CostCheck}{%
\begin{itemize}
  \item DS-V4.1-Flash：账户余额 ¥11.32 → ¥9.37，实扣 ¥1.95；按美元价目折算 \$0.39（人民币账户按人民币价目扣费）。
  \item DS-V4-Pro：按 token 计 \$3.33（全部在谷时）。运行中两次充值金额未记录，无法用余额核对。13 次调用因余额不足立即返回 402、不生成内容，充值后补跑。
  \item Kimi K3.1 Pro：三次充值共 \$17.50，最终余额 $-$\$0.09，实扣 \$17.59。成功返回的 35 次调用按 token 计 \$8.57。差额 \$9.02（按余额变化推算）对应 7 个未返回结果的请求：4 个在余额耗尽后挂起 1--1.5 小时，3 个被手动中断。
\end{itemize}}

\newcommand{\Analysis}{%
\begin{itemize}
  \item 共同结果：28 题三个模型都通过；\texttt{abc397\_d}、\texttt{arc195\_c} 三个模型都失败。只有 Kimi 通过：\texttt{arc196\_b}；只有 DS-V4-Pro 通过：\texttt{abc398\_f}。
  \item Kimi 失败的 5 题：
  \begin{itemize}
    \item \texttt{abc397\_d}、\texttt{arc195\_c}：答案错误。
    \item \texttt{abc398\_g}：运行错误（IndexError）；思考 11.9 万 token、31 分钟，正常结束。
    \item \texttt{abc398\_f}：运行错误，原因是评测框架限制（见下）。
    \item \texttt{abc399\_e}：返回与题目无关的拒答（内容是一部电影的营销手法），输入 610 token 正常，判为无代码；疑似服务端错误，未重跑。
  \end{itemize}
  \item DS-V4.1-Flash 失败：\texttt{abc399\_c}、\texttt{abc397\_d}、\texttt{arc195\_c} 答案错误；\texttt{abc398\_f}、\texttt{arc196\_b} 运行错误。DS-V4-Pro 失败：\texttt{abc397\_d}、\texttt{arc195\_c} 答案错误；\texttt{arc196\_b} 运行错误。
  \item 评测框架限制：LCB 判题用 StringIO 捕获标准输出，不支持 \texttt{sys.stdout.buffer}。35 题中只有 \texttt{abc398\_f} 上的 DS-V4.1-Flash 和 Kimi 用了它，均报 AttributeError；DS-V4-Pro 用 \texttt{print}，通过。读入用 \texttt{sys.stdin.buffer} 不受影响。
  \item 推理量：每题输出 token 中位数 DS-V4.1-Flash 12k、Kimi 14k、DS-V4-Pro 38k；困难题平均 22k / 31k / 61k。DS-V4-Pro 用 Flash 2.8 倍的困难题 token，多做对 1 道困难题。
  \item 调用稳定性：DeepSeek 余额不足时立即返回 402，补跑即可，已答结果不丢。Kimi 非流式请求在余额耗尽后挂起 1--1.5 小时，既不返回也不报错；改为流式 + 5 分钟空闲超时后，最后 3 题单路完成（每题 20--31 分钟）。
\end{itemize}}

\newcommand{\PerfBaseline}{25.7}
\newcommand{\PerfBaselineRange}{23.7--28.8}
\newcommand{\PerfBaseRss}{12.1}
\newcommand{\PerfNoiseMedian}{2.0}
\newcommand{\PerfNoiseHigh}{6.4}
\newcommand{\PerfRuns}{11{,}124}
\newcommand{\PerfMeasurable}{13}
\newcommand{\PerfCommon}{13}
\newcommand{\PerfThresh}{0.1}
\newcommand{\PerfSmallList}{\texttt{3777}、\texttt{3793}、\texttt{3805}、\texttt{3809}、\texttt{abc397\_e}、\texttt{abc397\_f}、\texttt{abc398\_b}、\texttt{abc398\_c}、\texttt{abc398\_d}、\texttt{abc398\_f}、\texttt{abc398\_g}、\texttt{abc399\_e}、\texttt{abc399\_f}、\texttt{abc400\_c}、\texttt{abc400\_d}、\texttt{arc195\_a}、\texttt{arc195\_c}、\texttt{arc195\_d}、\texttt{arc195\_e}、\texttt{arc196\_a}、\texttt{arc196\_b}、\texttt{arc196\_c}}
\newcommand{\PerfSmallN}{22}
\newcommand{\PerfRows}{%
\texttt{abc399\_c} & 中 & AC & 2517 & 0.23 & 50 & 0.26 & 22 & \textbf{0.22} & 44 \\
\texttt{abc399\_d} & 中 & AC & 1758 & 0.15 & 38 & 0.28\,(2.3$\times$) & 48 & \textbf{0.12} & 48 \\
\texttt{3717} & 难 & LC & 830 & \textbf{1.00} & 37 & 1.40 & 41 & 1.21 & 37 \\
\texttt{3744} & 难 & LC & 2137 & 0.27\,(5.0$\times$) & 32 & 0.67\,(12.5$\times$) & 32 & \textbf{0.05} & 32 \\
\texttt{3765} & 难 & LC & 9 & \textbf{0.14} & 16 & 0.14 & 16 & 1.24\,(9.1$\times$) & 16 \\
\texttt{3784} & 难 & LC & 391 & 0.11\,(2.2$\times$) & 22 & 0.06 & 28 & \textbf{0.05} & 22 \\
\texttt{3801} & 难 & LC & 0 & \textbf{0.11} & 26 & 0.11 & 26 & 0.93\,(8.3$\times$) & 18 \\
\texttt{abc397\_d} & 难 & AC & 0 & 0.14 & 12 & 0.14 & 12 & \textbf{0.11} & 12 \\
\texttt{abc397\_g} & 难 & AC & 0 & \textbf{0.01} & 14 & 0.02\,(2.5$\times$) & 14 & 0.23\,(36.1$\times$) & 14 \\
\texttt{abc400\_e} & 难 & AC & 0 & 0.43 & 36 & 0.39 & 32 & \textbf{0.26} & 24 \\
\texttt{abc400\_g} & 难 & AC & 2773 & \textbf{0.16} & 32 & 0.27 & 40 & 1.45\,(9.3$\times$) & 40 \\
\texttt{arc195\_b} & 难 & AC & 37 & \textbf{0.15} & 14 & 1.30\,(8.5$\times$) & 91 & 0.96\,(6.3$\times$) & 119 \\
\texttt{arc196\_d} & 难 & AC & 80 & \textbf{0.06} & 16 & 0.13\,(2.1$\times$) & 16 & 0.15\,(2.5$\times$) & 18 \\
}
\newcommand{\PerfSummaryRows}{%
Kimi & \textbf{1.29} & 7 & 2 & 16.7 & 26 & 50 \\
DS-Flash & 2.00 & 0 & 5 & 29.7 & 28 & 91 \\
DS-Pro & 2.74 & 6 & 6 & 52.8 & 24 & 119 \\
}
\newcommand{\KimiGM}{1.29}
\newcommand{\KimiWins}{7}
\newcommand{\KimiSlowTwo}{2}
\newcommand{\DSFlashGM}{2.00}
\newcommand{\DSFlashWins}{0}
\newcommand{\DSFlashSlowTwo}{5}
\newcommand{\DSProGM}{2.74}
\newcommand{\DSProWins}{6}
\newcommand{\DSProSlowTwo}{6}
\newcommand{\StaticCommon}{33}
\newcommand{\StaticRows}{%
Kimi & 5 & 7 & buffer.read 20、read 3、buffer.readline 1 \\
DS-Flash & 5 & 10 & buffer.read 17、readline 4、buffer.readline 2、read 1 \\
DS-Pro & 4 & 11 & buffer.read 14、read 6、readline 3、buffer.readline 1 \\
}
\newcommand{\StaticIoRows}{%
Kimi & buffer.read 20、read 3、buffer.readline 1 \\
DS-Flash & buffer.read 17、readline 4、buffer.readline 2、read 1 \\
DS-Pro & buffer.read 14、read 6、readline 3、buffer.readline 1 \\
}

\newcommand{\QualIssueTable}{%
\begin{tabular}{@{}llrrrrrrr@{}}
\toprule
 & & & \multicolumn{6}{c}{Pylint} \\
\cmidrule(lr){4-9}
模型 & 难度 & 份数 & Pylint 合计 & C & R & W & E & 得分 \\
\midrule
\multicolumn{9}{@{}l}{\textit{全部答案}} \\
Kimi & 中等 & 10 & 3.90 & 3.50 & 0.40 & 0.00 & 0.00 & 7.87 \\
 & 困难 & 24 & 7.83 & 4.62 & 1.62 & 1.54 & 0.04 & 8.37 \\
 & 合计 & 34 & \textbf{6.68} & \textbf{4.29} & \textbf{1.26} & 1.09 & 0.03 & \textbf{8.22} \\\midrule
DS-Flash & 中等 & 10 & 7.10 & 6.20 & 0.80 & 0.10 & 0.00 & 7.45 \\
 & 困难 & 25 & 10.44 & 7.04 & 2.84 & 0.56 & 0.00 & 8.16 \\
 & 合计 & 35 & 9.49 & 6.80 & 2.26 & \textbf{0.43} & \textbf{0.00} & 7.96 \\\midrule
DS-Pro & 中等 & 10 & 6.10 & 5.40 & 0.60 & 0.10 & 0.00 & 7.39 \\
 & 困难 & 25 & 15.24 & 10.44 & 2.92 & 1.88 & 0.00 & 7.81 \\
 & 合计 & 35 & 12.63 & 9.00 & 2.26 & 1.37 & \textbf{0.00} & 7.69 \\
\midrule
\multicolumn{9}{@{}l}{\textit{仅通过的答案}} \\
Kimi & 中等 & 10 & 3.90 & 3.50 & 0.40 & 0.00 & 0.00 & 7.87 \\
 & 困难 & 23 & 7.96 & 4.70 & 1.65 & 1.57 & 0.04 & 8.35 \\
 & 合计 & 33 & \textbf{6.73} & \textbf{4.33} & \textbf{1.27} & 1.09 & 0.03 & \textbf{8.20} \\\midrule
DS-Flash & 中等 & 10 & 7.10 & 6.20 & 0.80 & 0.10 & 0.00 & 7.45 \\
 & 困难 & 25 & 10.44 & 7.04 & 2.84 & 0.56 & 0.00 & 8.16 \\
 & 合计 & 35 & 9.49 & 6.80 & 2.26 & \textbf{0.43} & \textbf{0.00} & 7.96 \\\midrule
DS-Pro & 中等 & 10 & 6.10 & 5.40 & 0.60 & 0.10 & 0.00 & 7.39 \\
 & 困难 & 25 & 15.24 & 10.44 & 2.92 & 1.88 & 0.00 & 7.81 \\
 & 合计 & 35 & 12.63 & 9.00 & 2.26 & 1.37 & \textbf{0.00} & 7.69 \\
\bottomrule
\end{tabular}}
\newcommand{\QualRadonTable}{%
\begin{tabular}{@{}llrrrrrrrrr@{}}
\toprule
模型 & 难度 & 份数 & CC & MI & Func & LOC & LLOC & SLOC & CMT & BLN \\
\midrule
\multicolumn{11}{@{}l}{\textit{全部答案}} \\
Kimi & 中等 & 10 & 8.5 & 65.3 & 1.2 & 33.6 & 24.8 & 24.8 & 1.7 & 7.2 \\
 & 困难 & 24 & 17.7 & 59.8 & 2.2 & 82.2 & 59.1 & 60.8 & 5.5 & 16.2 \\
 & 合计 & 34 & \textbf{15.0} & \textbf{61.4} & 1.9 & 67.9 & 49.0 & 50.2 & 4.4 & 13.5 \\\midrule
DS-Flash & 中等 & 10 & 11.8 & 56.4 & 1.2 & 42.6 & 32.7 & 32.0 & 0.9 & 9.7 \\
 & 困难 & 25 & 18.6 & 50.1 & 1.5 & 93.6 & 71.3 & 71.0 & 3.6 & 19.9 \\
 & 合计 & 35 & 16.7 & 51.9 & 1.4 & 79.0 & 60.3 & 59.8 & 2.8 & 17.0 \\\midrule
DS-Pro & 中等 & 10 & 12.2 & 56.2 & 1.3 & 37.0 & 30.9 & 30.3 & 0.5 & 6.2 \\
 & 困难 & 25 & 21.3 & 54.2 & 2.0 & 101.4 & 80.5 & 80.0 & 5.3 & 17.1 \\
 & 合计 & 35 & 18.7 & 54.8 & 1.8 & 83.0 & 66.3 & 65.8 & 3.9 & 14.0 \\
\midrule
\multicolumn{11}{@{}l}{\textit{仅通过的答案}} \\
Kimi & 中等 & 10 & 8.5 & 65.3 & 1.2 & 33.6 & 24.8 & 24.8 & 1.7 & 7.2 \\
 & 困难 & 23 & 17.9 & 59.9 & 2.3 & 83.3 & 59.9 & 61.7 & 5.7 & 16.3 \\
 & 合计 & 33 & \textbf{15.0} & \textbf{61.5} & 1.9 & 68.2 & 49.2 & 50.5 & 4.5 & 13.5 \\\midrule
DS-Flash & 中等 & 10 & 11.8 & 56.4 & 1.2 & 42.6 & 32.7 & 32.0 & 0.9 & 9.7 \\
 & 困难 & 25 & 18.6 & 50.1 & 1.5 & 93.6 & 71.3 & 71.0 & 3.6 & 19.9 \\
 & 合计 & 35 & 16.7 & 51.9 & 1.4 & 79.0 & 60.3 & 59.8 & 2.8 & 17.0 \\\midrule
DS-Pro & 中等 & 10 & 12.2 & 56.2 & 1.3 & 37.0 & 30.9 & 30.3 & 0.5 & 6.2 \\
 & 困难 & 25 & 21.3 & 54.2 & 2.0 & 101.4 & 80.5 & 80.0 & 5.3 & 17.1 \\
 & 合计 & 35 & 18.7 & 54.8 & 1.8 & 83.0 & 66.3 & 65.8 & 3.9 & 14.0 \\
\bottomrule
\end{tabular}}
\newcommand{\QualDensityTable}{%
\begin{tabular}{@{}llrrrrrrrr@{}}
\toprule
 & & & \multicolumn{5}{c}{Pylint} & \multicolumn{2}{c}{Radon} \\
\cmidrule(lr){4-8}\cmidrule(lr){9-10}
模型 & 难度 & 份数 & Pylint 合计 & C & R & W & E & CC & Func \\
\midrule
Kimi & 中等 & 10 & 15.73 & 14.11 & 1.61 & 0.00 & 0.00 & 34.27 & 4.84 \\
 & 困难 & 23 & 12.90 & 7.61 & 2.68 & 2.54 & 0.07 & 28.96 & 3.66 \\
 & 合计 & 33 & \textbf{13.32} & \textbf{8.58} & \textbf{2.52} & 2.16 & 0.06 & 29.75 & 3.84 \\\midrule
DS-Flash & 中等 & 10 & 22.19 & 19.38 & 2.50 & 0.31 & 0.00 & 36.88 & 3.75 \\
 & 困难 & 25 & 14.71 & 9.92 & 4.00 & 0.79 & 0.00 & 26.21 & 2.14 \\
 & 合计 & 35 & 15.85 & 11.37 & 3.77 & \textbf{0.72} & \textbf{0.00} & \textbf{27.84} & 2.39 \\\midrule
DS-Pro & 中等 & 10 & 20.13 & 17.82 & 1.98 & 0.33 & 0.00 & 40.26 & 4.29 \\
 & 困难 & 25 & 19.04 & 13.04 & 3.65 & 2.35 & 0.00 & 26.64 & 2.45 \\
 & 合计 & 35 & 19.18 & 13.67 & 3.43 & 2.08 & \textbf{0.00} & 28.43 & 2.69 \\
\bottomrule
\end{tabular}}
\newcommand{\QualSonar}{0}
\newcommand{\QualBullets}{%
  \item Pylint（Kimi / DS-Flash / DS-Pro，每份）：问题 6.73 / 9.49 / 12.63，得分 8.20 / 7.96 / 7.69。DS-Pro 最多，主要是命名不合规（C0103，共 25 / 110 / 174 条）。
  \item 警告（W）Kimi 每份 1.09，高于 DS-Flash 的 0.43，但集中在 2 份答案的循环内闭包（W0640 33 条）；DS-Pro 的 W0640 有 34 条，集中在 1 份。未使用的变量（W0612）：Kimi 0、DS-Flash 3、DS-Pro 8 条。
  \item Radon（每份）：Kimi 代码最短（SLOC 50.5 / 59.8 / 65.8）、可维护性指数最高（61.5 / 51.9 / 54.8）、圈复杂度最低（15.0 / 16.7 / 18.7）。
  \item 按每 100 SLOC：Pylint 问题 13.32 / 15.85 / 19.18，排序不变；但 Kimi 的圈复杂度密度（29.75 / 27.84 / 28.43）和函数密度（3.84 / 2.39 / 2.69）最高，即代码更短、分支和函数更密。}


\begin{document}

\begin{center}
{\LARGE\bfseries LiveCodeBench 代码能力评测}\\[0.4em]
{\large Kimi K3.1 Pro / DeepSeek-V4.1-Flash / DeepSeek-V4-Pro}\\[0.3em]
{\small 2026-10-23（UTC） \quad LiveCodeBench release\_v6 最新 35 题（中等 10、困难 25）}
\end{center}

\section*{结论}
\begin{itemize}
  \item 正确率：\KeySummaryPass
  \item 速度：\KeySummarySpeed
  \item 费用：\KeySummaryCost
  \item 代码：\KeySummaryCode
  \item \KeySummaryTakeaway
\end{itemize}

\section{模型}

\begin{center}
\small
\begin{tabularx}{\linewidth}{@{}>{\raggedright}p{2.2cm}*{3}{>{\raggedright\arraybackslash}X}@{}}
\toprule
 & Kimi K3.1 Pro & DeepSeek-V4.1-Flash & DeepSeek-V4-Pro \\
\midrule
API 模型名 & \texttt{kimi-k3.1-pro} & \texttt{deepseek-flash} & \texttt{deepseek-v4-pro} \\
最近更新 & K3.1 系列首发 & 2026-10-17（1017 版） & 2026-08-13（GA） \\
规模 & 3T 级，K3（2.8T）后训练 & 552B MoE，激活 8B 输入 / 16B 输出 & 未公布 \\
权重 & 开放 & 开放，MIT & — \\
上下文 / 最大输出 & 1M / 默认 131,072 & 1M / 393,216 & 1M / 393,216 \\
输入 & 文本、图像 & 文本、图像 & 文本 \\
思考档位 & 恒开；low/high/max，默认 max & low/high/max，默认 high & low/high/max，默认 high \\
采样 & 固定 T=1.0、top\_p=0.95 & 推荐 T=1.0 & — \\
单价（美元/百万 token，输入 / 缓存 / 输出） & 2 / 0.2 / 10，无峰谷 & 0.3 / 0.006 / 1.2，谷时半价 & 1.32 / 0.044 / 3.96，谷时半价 \\
\bottomrule
\end{tabularx}
\end{center}

\subsection{Kimi K3.1 Pro}
\begin{itemize}
  \item 变更：在 K3 上大幅增加后训练算力，覆盖数百个领域联合 RL；重做 RL 框架、长上下文推理、沙箱、rollout 调度。
  \item 官方提升方向：编程、3D 建模（直接写 Blender / CAD 脚本）、金融研究（内部 Financial Research 比 K3 高 11.6 分）、智能体。
  \item 安全：中期训练与 RL 数据剔除网络安全内容；xdev-bench-hard pass@3 81.0\%（不剔除 90.5\%）。
  \item 高速版 \texttt{kimi-k3.1-pro-highspeed}：速度 2--3 倍，价格 2 倍。K2 系列 2026 年底下线。
  \item 使用限制：充值 $\ge$1 美元解锁；累计充值决定并发 / TPM；temperature、top\_p、n、penalty 不可传。
  \item 官方文档中的编程基准分数为占位符（“[xx.x\%]”），无可引用数据。
\end{itemize}

\subsection{DeepSeek-V4.1-Flash}
\begin{itemize}
  \item 2026-09-10 发布：20 层因果编码器 + 20 层解码器；MoE 1 共享 + 384 路由专家，每 token 激活 6 个；KV 缓存 890 字节/token，HBM 为上一代 1/4、SSD 为 1/8；原生视觉；降价；V4-Flash 下线。
  \item 2026-10-17 更新为 1017 版（新 / 旧）：Terminal-Bench 4.0 42.3 / 31.2，DeepSWE v1.1 77.0 / 74.2，ProgramBench 29.1 / 20.3，NL2Repo 68.2 / 65.4，HLE 42.6 / 36.8；输出更短、耗时更少；价格与调用方式不变。
  \item Codeforces 3471（max 档；同表 V4-Pro 3348）。输出单价为 V4-Pro 的 30\%、Kimi K3.1 Pro 的 12\%。
\end{itemize}

\subsection{DeepSeek-V4-Pro}
\begin{itemize}
  \item 2026-04-24 首发；2026-08-13 GA：HLE 42.7 / 60.0（无 / 有工具），Terminal-Bench 2.1 87.9，DeepSWE 62.7，NL2Repo 61.5，Toolathlon-Verified 74.1；引入 low/high/max 三档；2026-08-16 起峰谷定价。
  \item V4.1-Flash 新闻稿称 9 月 14 日起该模型名改由 V4.1-Flash 承接，更新日志称继续提供、计费不变。本次 API 返回模型名 \DSProServed，按 V4-Pro 单价扣费。
  \item 仅文本输入；单价为 V4.1-Flash 的 3.3 倍；DeepSeek 自报多数基准低于 V4.1-Flash。
\end{itemize}

\section{评测方法}
\begin{center}
\small
\begin{tabularx}{\linewidth}{@{}>{\raggedright}p{2.4cm}X@{}}
\toprule
项目 & 设置 \\
\midrule
评测集 & LiveCodeBench release\_v6（ICLR 2025），官方提示词与官方测试数据，全部测试通过才算通过 \\
指标 & pass@1，每题 1 次 \\
调用 & 官方 API；\texttt{BAI-AI-Lab/codingBench}，\texttt{cb livecodebench} \\
思考档位 & 各模型默认（Kimi max，DeepSeek high） \\
采样 & DeepSeek：temperature 0.2（LCB 官方）；Kimi：固定 1.0 \\
输出上限 & DeepSeek 262,144；Kimi 131,072（默认）；\TruncNote \\
并发 & V4.1-Flash 16，V4-Pro 8，Kimi 1--4；延迟含服务端排队 \\
计量 & 每次调用记录起止时间与服务端返回的 token（含思考） \\
费用 & 官方单价 $\times$ token，DeepSeek 按调用时段区分峰谷 \\
判题 & AtCoder：每份答案作为独立 Python 进程，用真实标准输入输出运行全部官方测试（每个测试 6 秒），逐 token 比对；\texttt{abc397\_d}、\texttt{arc195\_c} 用按题意编写的校验器（前者第 41 个测试的标准答案不满足题意，后者允许任意合法构造）。LeetCode：LCB 的函数调用判题 \\
\bottomrule
\end{tabularx}
\end{center}

\section{评测题目}

\subsection{分布}
release\_v6 最新 100 题中，按难度各取最新：中等 10、困难 25；发布于 2025-03-15 至 2025-04-06。

\begin{center}
\small
\begin{tabular}{@{}lrrr@{}}
\toprule
难度 & AtCoder & LeetCode & 合计 \\
\midrule
中等 & 7 & 3 & 10 \\
困难 & 19 & 6 & 25 \\
合计 & 26 & 9 & 35 \\
\bottomrule
\end{tabular}
\qquad
\begin{tabular}{@{}lrl@{}}
\toprule
来源 & 题数 & 形式 \\
\midrule
AtCoder ABC 397--400 & 17 & 完整程序，标准输入输出 \\
AtCoder ARC 195--196 & 9 & 完整程序，标准输入输出 \\
LeetCode 周赛 / 双周赛 & 9 & 补全给定 Python 函数 \\
\bottomrule
\end{tabular}
\end{center}

\subsection{示例}
\begin{center}
\small
\begin{tabularx}{\linewidth}{@{}>{\raggedright}p{2.6cm}X@{}}
\toprule
题目 & 内容 \\
\midrule
\textbf{ABC 400 C}\newline “2\^{}a b\^{}2”\newline 中等，2025-04-05 &
\textbf{题意}：$X = 2^a b^2$（$a,b \ge 1$）为好数；给定 $N \le 10^{18}$，求 $[1,N]$ 内好数个数。$N=20$ 时为 5（2, 4, 8, 16, 18）。\newline
\textbf{考点}：好数恰为 $2m^2$ 或 $4m^2$ 两类且互不重叠，答案 $\lfloor\sqrt{N/2}\rfloor+\lfloor\sqrt{N/4}\rfloor$；须用整数开方，浮点 \texttt{sqrt} 在 $10^{18}$ 附近出错。\newline
\textbf{结果}：\ExampleAResult \\
\addlinespace
\textbf{ARC 195 D}\newline “Swap and Erase”\newline 困难，2025-03-23 &
\textbf{题意}：可交换相邻元素，或删除开头一段全部相等的前缀；求删空序列的最少操作数。$T \le 10^5$，$\sum N \le 2\times 10^5$。\newline
\textbf{考点}：交换与删除的组合优化，须近线性时间。样例：$(1,1,2,1,2)$ 为 3（1 次交换 + 2 次删除），$(4,2,1,3)$ 为 4（不交换）。\newline
\textbf{结果}：\ExampleBResult \\
\addlinespace
\textbf{LeetCode 3801}\newline “Count Beautiful Numbers”\newline 困难，2025-03-15 &
\textbf{题意}：各位数字之积能被各位数字之和整除的正整数为美丽数；给定 $1 \le l \le r < 10^9$，求区间内个数。补全 \texttt{beautifulNumbers(self, l, r)}。\newline
\textbf{考点}：数位 DP；含 0 时积为 0 必满足，否则积为 $2^a3^b5^c7^d$，状态有限；暴力枚举超时。\newline
\textbf{结果}：\ExampleCResult \\
\bottomrule
\end{tabularx}
\end{center}

\section{评测结果}

\subsection{总表}
\begin{center}
\small
\begin{tabular}{@{}lccc@{}}
\toprule
指标 & Kimi K3.1 Pro & DS-V4.1-Flash & DS-V4-Pro \\
\midrule
\OverallRows
\bottomrule
\end{tabular}
\end{center}
{\footnotesize 粗体为该行最优。1 题 = 2.9 个百分点。有效运行时间 = 成功请求覆盖的时长。}

\begin{center}
\begin{tikzpicture}
\begin{axis}[ybar=2pt, width=0.86\linewidth, height=5.4cm, bar width=16pt, area legend, legend image post style={draw=none},
  ymin=0, ymax=114, ytick={0,20,40,60,80,100}, ylabel={pass@1（\%）}, symbolic x coords={中等,困难,总计}, xtick=data, enlarge x limits=0.25,
  nodes near coords={\pgfmathprintnumber[fixed,precision=0]{\pgfplotspointmeta}}, nodes near coords style={font=\scriptsize},
  legend style={at={(0.5,1.03)},anchor=south,legend columns=3,draw=none,font=\small,/tikz/every even column/.append style={column sep=10pt}},
  ymajorgrids, grid style={black!10}]
\addplot[fill=kimi,draw=none] coordinates {\KimiCoords};
\addplot[fill=dsflash,draw=none] coordinates {\DSFlashCoords};
\addplot[fill=dspro,draw=none] coordinates {\DSProCoords};
\legend{Kimi K3.1 Pro, DeepSeek-V4.1-Flash, DeepSeek-V4-Pro}
\end{axis}
\end{tikzpicture}
\end{center}

\subsection{按难度}
\begin{center}
\small
\begin{tabular}{@{}llcccc@{}}
\toprule
难度 & 模型 & 通过 & 延迟中位数（秒） & 平均输出 token & 最多输出 token \\
\midrule
\DiffRows
\bottomrule
\end{tabular}
\end{center}

\subsection{分析}
\Analysis

\subsection{逐题结果}
{\footnotesize 每个模型：结果 / 输出 token（含思考）/ 延迟（秒）。$^{\dagger}$ 评测中 API 返回的内容与题目无关，人工另行调用验证正确，按通过计。}
{\scriptsize
\setlength{\tabcolsep}{3pt}
\begin{longtable}{@{}llclrrlrrlrr@{}}
\toprule
题号 & 题目 & 难度 & \multicolumn{3}{c}{Kimi K3.1 Pro} & \multicolumn{3}{c}{DS-V4.1-Flash} & \multicolumn{3}{c}{DS-V4-Pro} \\
\midrule
\endfirsthead
\toprule
题号 & 题目 & 难度 & \multicolumn{3}{c}{Kimi K3.1 Pro} & \multicolumn{3}{c}{DS-V4.1-Flash} & \multicolumn{3}{c}{DS-V4-Pro} \\
\midrule
\endhead
\bottomrule
\endlastfoot
\ProblemRows
\end{longtable}}

\section{代码性能与静态指标}\label{sec:perf}

\subsection{测量方法}
\begin{center}
\small
\begin{tabularx}{\linewidth}{@{}>{\raggedright}p{2.4cm}X@{}}
\toprule
项目 & 设置 \\
\midrule
对象 & 通过的答案（Kimi 33 份，不含 \texttt{abc398\_g}、\texttt{abc399\_e}；DeepSeek 各 35 份） \\
运行 & 每次运行都是一个全新的 Python 3.11 进程（LCB 环境），绑定到单个 CPU 核；共 \PerfRuns{} 次，输出全部正确 \\
计时 & AtCoder：进程 CPU 时间（rusage）扣除启动基线 \PerfBaseline{} ms；LeetCode：只计 \texttt{Solution} 方法调用的 CPU 时间 \\
内存 & 进程峰值常驻内存（ru\_maxrss）；空程序为 \PerfBaseRss{} MB \\
重复 & 每个测试 3 次、三个模型轮流先跑，取中位数；同一测试 3 次的离散度中位数 \PerfNoiseMedian\%，P90 \PerfNoiseHigh\% \\
可比题 & 最慢答案扣除启动后仍 $\ge$ \PerfThresh{} 秒的题，共 \PerfMeasurable{} 题；其余 \PerfSmallN{} 题官方测试太小，耗时主要是启动开销，不比较 \\
\bottomrule
\end{tabularx}
\end{center}

\subsection{逐题性能}
\begin{center}
\small
\setlength{\tabcolsep}{4pt}
\begin{tabular}{@{}lccr rr rr rr@{}}
\toprule
题号 & 难度 & 类型 & 最大输入 KB & \multicolumn{2}{c}{Kimi K3.1 Pro} & \multicolumn{2}{c}{DS-V4.1-Flash} & \multicolumn{2}{c}{DS-V4-Pro} \\
\cmidrule(lr){5-6}\cmidrule(lr){7-8}\cmidrule(l){9-10}
 & & & & CPU 秒 & MB & CPU 秒 & MB & CPU 秒 & MB \\
\midrule
\PerfRows
\bottomrule
\end{tabular}
\end{center}
{\footnotesize CPU 秒为最慢一个测试的 CPU 时间（3 次取中位数，扣除启动）；粗体为该题最快；括号内为相对最快者的倍数（$\ge 2$ 时标出）。MB 为该题各测试中的最大峰值内存。AC = AtCoder，LC = LeetCode。}

\subsection{汇总}
\begin{center}
\small
\begin{tabular}{@{}lcccccc@{}}
\toprule
模型 & 相对最快倍数 & 最快 & 慢 $\ge 2$ 倍 & CPU 合计 & \multicolumn{2}{c}{峰值内存（MB）} \\
 & （几何平均） & （题） & （题） & （秒） & 中位数 & 最大 \\
\midrule
\PerfSummaryRows
\bottomrule
\end{tabular}
\end{center}
{\footnotesize 基于 \PerfCommon{} 道可比题（三个模型都通过）。CPU 合计为这些题所有测试的 CPU 时间之和。“最快”按严格最小值计；有 3 题第二名与最快相差不到 5\%（噪声范围内）：\texttt{3765}、\texttt{3801} 上 Flash 比 Kimi 慢 3\% 和 2\%，\texttt{abc399\_c} 上 Kimi 比 Pro 慢 3\%。按并列计，Flash 的最快题数为 2。}

\begin{itemize}
  \item \texttt{arc195\_b}（Flash 慢 8.5 倍，内存 91 MB）：Flash 枚举 A、B 所有不同取值的两两组合，复杂度 $O(U\cdot V)$，并把重复的和逐个展开存入列表；Kimi 先按通配数量选出必须参与配对的候选值，再用二分查找限定候选和，0.15 秒。
  \item \texttt{abc400\_g}（Pro 慢 9.3 倍）：Pro 为每组数据建图跑最小费用流；Kimi 用“取最大值后只修正奇偶”的贪心，每类只保留前两个候选，$O(N\log N)$。
  \item \texttt{3765}（Pro 慢 9.1 倍）：Pro 按子数组个数分层、每层做凸包优化，$O(n^2\log n)$；Kimi 直接用 $O(n^2)$ 的后缀 DP。
  \item \texttt{3801}（Pro 慢 8.3 倍）：Pro 对每个可能的数位和 $S$ 单独建 [剩余位数][和][积 mod $S$] 计数表；Kimi 用以（位置，和，积）为状态的记忆化数位 DP。
  \item \texttt{abc397\_g}：Pro 0.23 秒、Kimi 0.01 秒，倍数大但绝对值都很小。
  \item 全部可比题的最慢测试都在 1.5 秒以内，未接近 AtCoder 常见的 2 秒时限；但评测机器与 AtCoder 不同，不能据此判断能否在原平台通过。
\end{itemize}

\subsection{静态代码质量}\label{sec:static}
按 Huang 等（arXiv:2610.19167，ICECCS 2026）的口径：Pylint 3.3.5（默认配置）、Radon 6.0.1（命令行）\ifnum\QualSonar=1、SonarQube 10.4 Community（本机部署，默认质量配置）\fi；每份答案原样分析，不补 import；不合成总分。Kimi 的 \texttt{abc399\_e} 没有代码，不计入；“全部答案”与“仅通过”只差 Kimi 的 \texttt{abc398\_g} 一份。

\begin{center}
\footnotesize
\setlength{\tabcolsep}{3.5pt}
\QualIssueTable
\end{center}
{\footnotesize 每份答案平均问题数。C 规范、R 重构、W 警告、E 错误；得分 0--10，越高越好\ifnum\QualSonar=1；SonarQube 为可维护性、可靠性、安全三类问题数\fi。粗体为合计行中最优。每份答案都有一条“缺少末尾换行”（C0304），来自答案提取，三个模型相同。Kimi 的 1 个 E 是 \texttt{arc196\_c} 导入 numpy（分析环境未安装 numpy）。}

\begin{center}
\footnotesize
\setlength{\tabcolsep}{4pt}
\QualRadonTable
\end{center}
{\footnotesize 每份答案平均。CC 为 \texttt{radon cc -j} 全部条目之和（类与它的方法各计一条）；MI 为可维护性指数，0--100，越高越好；Func 为条目数；LOC、LLOC、SLOC、CMT、BLN 分别为总行、逻辑行、有效行、注释行、空行。}

\begin{center}
\footnotesize
\setlength{\tabcolsep}{4pt}
\QualDensityTable
\end{center}
{\footnotesize 仅通过的答案，每 100 有效行（SLOC）的数量：各题之和 $\div$ SLOC 之和 $\times$ 100。}

\begin{itemize}
\QualBullets
\end{itemize}

{\footnotesize 补充（不属于上述口径）：}
\begin{center}
\small
\begin{tabular}{@{}lccl@{}}
\toprule
模型 & 嵌套深度中位数 & 嵌套深度最大 & AtCoder 读入方式（题数） \\
\midrule
\StaticRows
\bottomrule
\end{tabular}
\end{center}
{\footnotesize 基于 \StaticCommon{} 道三个模型都通过的题。嵌套深度为语法树中复合语句的最大嵌套层数。Kimi 采样温度 1.0、DeepSeek 0.2，风格差异部分来自采样。}

\section{局限}
\begin{itemize}
  \item 35 题、每题 1 次：相差 1--2 题不显著；两个 DeepSeek 模型均为满分，这批题已不足以区分二者。
  \item 采样温度不同：DeepSeek 0.2，Kimi 固定 1.0。
  \item 延迟受并发和服务端负载影响。
  \item 题目早于三个模型发布，不排除训练数据污染。
  \item 成绩依赖本评测的判题程序和校验器。
  \item Kimi \texttt{abc399\_e} 按人工另行调用的结果计为通过，未在本评测流程中复现。
\end{itemize}

\section*{来源}
{\small
\begin{itemize}
  \item Kimi：\url{https://platform.kimi.ai/docs/guide/kimi-k3-1-pro-quickstart}，\url{https://platform.kimi.ai/docs/models}，\url{https://platform.kimi.ai/docs/pricing/chat}
  \item DeepSeek：\url{https://api-docs.deepseek.com/updates/}，\url{https://www.deepseek.com/en/news/deepseek-v4-1-flash/}，\url{https://huggingface.co/deepseek-ai/DeepSeek-V4.1-Flash}，\url{https://api-docs.deepseek.com/quick_start/pricing}
  \item 上下文、最大输出、默认档位：两家 API \texttt{/v1/models}（2026-10-23）
  \item LiveCodeBench：\url{https://livecodebench.github.io}
\end{itemize}}

\end{document}
